\documentclass[12pt,a4paper]{article}
\usepackage[UTF8]{ctex}
\usepackage[margin=2.5cm]{geometry}
\usepackage{hyperref}
\usepackage{setspace}
\usepackage{graphicx}
\usepackage{amsmath}
\usepackage{enumitem}
\usepackage{natbib}

\setstretch{1.5}
\setlength{\parindent}{2em}
\setlength{\parskip}{0pt}

\begin{document}

\title{\textbf{开题报告：基于MRAG-Bench的多模态检索器性能评估与优化}}
\author{}
\date{}
\maketitle

\section*{一、研究目的}

实现一种新的多模态检索器以提高多模态大语言模型的多模态检索能力。

\section*{二、本课题的目的及研究意义}

\subsection*{2.1 研究目的}

近年来，以Transformer\cite{vaswani2017attention}及其变体为原理的大规模语言模型（LLM）在许多语言相关的任务中取得了非凡的成功。通过在大量高质量的指令数据集上进行预训练，LLMs可以学习各种各样的语言模式、结构和事实知识，展现出了强大的自然语言理解与自然语言生成的能力\cite{minaee2024llm}。LLM是新一轮科技革命的开端。2023年ChatGPT的发布宣布了AI时代的到来，2024年闭源开源大模型如雨后春笋纷纷破土而出，OpenAI o1完成基于强化学习（RL）做推理后，Deepseek用强化学习将思维链（CoT）低成本普及开来\cite{wei2022cot}。2025年，为了解决大模型知识落后和幻觉问题，RAG\cite{lewis2020rag}应运而生，诞生了Dify\cite{dify2023}等一众知识增强框架。Manus\cite{shen2025manus}的爆火带动了智能体（Agent）\cite{wang2023agent}概念的出现。随后Claude引领了Agent的浪潮，Claude Code\cite{anthropic2024claudecode}、MCP\cite{anthropic2024mcp}、Skills\cite{langchain2024skills}、Clawdbot\cite{clawdbot2025}等技术涌现赋予了Agent调用工具并系统性地与虚拟环境和真实计算机系统互动的能力。目前各家大模型公司已经在编程、解题、医学顾问、语言学习等领域相互刷榜，以基座模型赋能众多垂直领域。

但大模型在真实性与可用性上同时存在着信息时效性（outdated/frozen knowledge）与幻觉（hallucination）等问题\cite{lewis2020rag,gao2023rag}。虽然这些基座模型通常在多样且广泛的数据库上进行训练，但这些数据库可能无法涵盖高度专业领域或实时信息更新所需的深度知识，掣肘了LLM在需要高精度和最新知识领域的性能。这就是信息时效性问题。当知识与信息完备时，LLM的推理能力能够赋予其创新能力，产生全新的见解，但当查询的内容超出了其训练的知识范围与其掌握的落实信息时，LLM可能会产生基于模式而非严格推理所得出的可能导致误导、不准确甚至完全捏造的思考与答案，这种现象称为幻觉。

RAG是一种允许LLMs在推理阶段动态集成最新信息的技术\cite{lewis2020rag,gao2023rag}。通过在生成前增加retrieve步骤，从外部知识库中查询相关的实时信息，利用外部文本知识，将其与查询一同输入LLM进行问答，从而使模型的回答以检索到的事实为依据\cite{gao2023rag}。这一机制让模型的知识库能够得到动态更新，而不局限于训练时落后的静态语料，同时其可追溯性也大幅度降低了编造信息的概率。可以说，RAG将LLM的问答逻辑从闭卷考试转换成了开卷考试的范式，显著增强了回答的准确性与时效性。

多模态信息指的是包含文本、图片、视频、音频等不同形式的数据格式。最初的LLM都聚焦在纯文本这一单一模态的训练与生成，而语言并非建模真实世界、满足用户需求的唯一媒介。许多研究者开始诉诸基于多模态的数据训练多模态大语言模型（Multimodal LLMs/LVLMs/MLLMs），把"语言生成能力"扩展到图像、视频、文档、表格与语音等多模态输入输出，来实现多模态数据理解与回答生成的需求\cite{mei2025multimodalrag,abootorabi2025ask}。这种方法不仅提高了问答系统的质量，而且通过将响应与事实、多模态知识联系起来，大大降低了幻觉的发生率。

在多模态的需求与技术赋能下，信息时效性与幻觉等问题有了新的形式与挑战。一方面，多模态模型在视觉理解与生成中会出现"凭空描述不存在物体/属性"的现象，学界已专门构建了面向视觉幻觉的评测基准，例如POPE用于评估多模态模型的对象幻觉问题\cite{li2023pope}。另一方面，多模态模型的知识主要来自预训练语料与参数记忆，面对快速变化、长尾、领域私有或实时更新的信息时，很容易因为"知识冻结"而出现过时回答\cite{gao2023rag}。

因此，RAG逐渐走向mRAG，用于管理外挂与基座多模态大模型的多模态知识库。多模态大语言模型与多模态检索增强生成（Multimodal RAG，mRAG/MM-RAG/MRAG）应运而生，逐渐成为研究的热点\cite{mei2025multimodalrag,abootorabi2025ask}。mRAG将传统RAG的纯文本范式框架拓展到了支持图像、视频等多种模态的外部知识的范式，在检索阶段融合跨模态的外部知识，并在生成阶段使用多模态大模型的基座来产生含有多模态数据的回答\cite{abootorabi2025ask}。通过引入视觉信息，mRAG系统能够提供更全面、直观、动态的回答。

学术上，mRAG把"信息检索"与"多模态生成/推理"耦合成端到端系统，迫使研究者回答一系列关键问题：如何表达多模态查询、如何对齐跨模态语义空间、如何评估"检索正确但生成仍错"的链路疾病、如何在噪声检索下保证稳健与可信\cite{abootorabi2025ask}。多模态RAG的输入/输出组合空间极大，统计表明常用模态输入输出总计54种组合，但已有研究只覆盖其中18种，表明这一方向仍处在快速发展且大量空白待填的阶段\cite{abootorabi2025ask}。

因此，本课题的目的在于：基于多模态检索器实现多模态检索的效果提升。而本课题的研究内容就是为了提高多模态检索，在根据查询query调用检索器获取与查询相关的图片的环节，提升检索器检索到有效信息（GT）的准确率。

\subsection*{2.2 研究意义}

优化检索器与Diffusion模型驱动训练的mRAG如何赋能具身智能、脑机接口、类脑智能与AGI。

\textbf{(1) 具身智能}

在具身智能领域，具身智能（机器人）面对的是强多模态、强相关、强时空连续数据：视觉、力觉、位姿、地图、任务步骤与语言指令相互耦合。Embodied-RAG明确指出：机器人可以无限探索与学习，但知识必须"可搜索且可行动"；同时，传统RAG技术并不能直接迁移到具身领域，因为具身数据多模态且高度相关、感知需要抽象\cite{xie2024embodiedrag}。为此它提出在具身基础模型上叠加"非参数记忆系统"，并以语义森林（semantic forest）组织不同分辨率的语言描述，实现跨平台、跨环境的解释与导航查询。

通过对意图对齐、长文本支持、细粒度证据检索、规划式检索，结合扩散模型生成/反馈信号提升检索器与生成模型协同等一系列的面向mRAG的检索器优化的核心意义在于把多模态大模型从"会说/会看但不可靠"升级为"能检索、能对齐证据、能在动态世界中持续更新并可解释地行动"。这一点对具身智能、脑机接口、类脑智能与AGI（通用智能体）是共通且关键的基础设施级能力。

具身场景的检索需求往往不是"找最相似的图/句子"，而是"找与当前任务意图相关的经验片段"（类似场景的失败案例、某类物体的可供性、某种操作步骤的关键例外）。当检索器仍停留在相似度最近邻范式时，容易把"看起来像但行动不适用"的证据检索回来，导致规划与控制出现错误甚至安全风险。我们的研究将把mRAG从"给机器人加个向量库"提升为"机器人可检索的长期记忆系统"：提高长任务成功率、降低错误规划与幻觉式行动、增强对长尾场景适应，并为人机协作提供"为什么这么做"的证据链与可解释性。ELLMER的研究表明，通过将RAG基础设施与具身智能结合，可以在不可预测环境下使机器人完成长视野任务并适应条件变化\cite{mon2024ellmer}。

\textbf{(2) 脑机接口}

BCI尤其是"脑控输入法/辅助沟通"对可靠性要求极高。Carìa等人（2025）的研究指出：LLM会产生看似连贯但错误的输出（hallucinations），在BCI沟通中必须缓解；并明确提出RAG（检索增强生成）作为有前景策略，通过融合外部知识库减少幻觉、提升事实准确性，从而提高沟通可靠性\cite{caria2025bci}。该文还指出BCI长期应用需要持续学习，而深度模型存在灾难性遗忘，LoRA/Adapter与RAG可部分缓解，但仍缺乏完善方案。

BCI场景的"正确输出"往往依赖个体化上下文（用户常用词、专业领域词汇、对话上下文、近期事件），这类信息动态变化且高度私域。通用相似度检索器难以做到"个体意图+私域语境"的精确召回，检索噪声会直接变成错误建议或误沟通。Memory Recall框架提出了一种检索增强的脑解码重建方法，通过检索视觉经验来补全不完整的脑信号，为BCI中的多模态检索提供了新的思路\cite{zhao2025memoryrecall}。

\textbf{(3) 类脑智能}

类脑智能关注"情景记忆（episodic）与语义记忆（semantic）如何协同"。GENESIS直接把情景记忆建模为RAG架构：海马体系统在检索增强生成框架内进行情景编码与检索，从而支持回忆、序列回忆效应、要点化（gist）与构造性想象等现象\cite{genesis2025}。这提供了一个非常强的"研究意义锚点"：RAG/mRAG与类脑记忆机制在形式上同构。

类脑系统需要的不只是"召回相似片段"，而是"按目标重组经验"，并在不确定时检索更多证据来抑制幻想。纯相似度检索很难实现"目标驱动的经验选择、对抗性干扰下的稳定回放、跨模态线索的记忆补全"。

\section*{三、本课题的国内外研究现状}

\subsection*{3.1 多模态检索增强生成（mRAG）的框架与关键问题}

mRAG领域的研究目前聚焦两个阶段与一系列不同的子问题。典型的mRAG系统可以分为检索与生成两个阶段。首先，在检索阶段，系统会根据用户的文本、图片、音频等多模态查询信息，从大模型知识库中检索相关的外部信息，包括文本文档、图片、表格、知识图谱、甚至是视频片段与字幕。然后将检索到的结果与原始查询信息一并输入多模态大模型中，以此生成融合了检索知识的答案\cite{mei2025multimodalrag,abootorabi2025ask}。

\textbf{挑战一：降低幻觉与保证时效。}mRAG通过引入外部知识有效地缓解了幻觉问题，这是该领域研究的核心动机之一\cite{gao2023rag,li2023pope}。多模态LLM在回答过程中可以引用检索到的事实性图文避免凭空捏造不确定的内容。对于需要视觉佐证的问题，mRAG提供的图像可以作为模型回答的依据，提升了回答的可信度。同时，由于检索库可以不断更新最新的信息，所以mRAG模型的回答能够反映最新的事实和动态\cite{mei2025multimodalrag}。

模型不仅会"编造文本事实"，还会"编造视觉事实"（例如描述图像中不存在的对象或属性）。在高风险领域（如医疗辅助诊断、自动驾驶、脑机接口等），此类幻觉会直接侵蚀系统可信度与安全性。信息时效性问题同样突出。LLM/MLLM主要依赖静态训练数据，容易产生"过时回答"或在缺少最新外部信息时进行猜测\cite{gao2023rag}。

\textbf{挑战二：检索策略与框架演化。}早期的检索增强生成往往为每个查询固定检索若干条文本。但在多模态场景下，不同查询对外部信息的需求差异更大，盲目检索所有模态既低效又可能引入噪音。为此，近期工作开始探索自适应的检索规划。例如，Windsock框架引入查询依赖的决策模块，智能判断是否需要检索，以及需要检索哪些模态\cite{zhao2025windsock}。这种按需检索的策略减少了不必要的开销，据报告可在保持回答质量的同时将检索次数减少约40\%。这些改进代表了mRAG框架从静态到动态、从单轮到多轮的演进趋势，使系统更高效且鲁棒。

\textbf{挑战三：多模态对齐与表示学习。}高效的跨模态检索依赖于良好的多模态表示对齐技术，即将文本、图像等映射到统一向量空间以比较相关性\cite{radford2021clip}。为此，大量研究借助对比学习预训练模型如CLIP和ALIGN，将图文特征投影到同一语义空间来计算相似度\cite{radford2021clip,gomez2021align}。不过，基本的CLIP存在细粒度理解不足和文本长度受限等问题\cite{xie2025fgclip,zhang2024longclip}。针对这些瓶颈，业界提出了多种改进的通用多模态检索模型\cite{xie2025fgclip,zhang2024longclip,ma2022eiclip}。

\subsection*{3.2 通用多模态检索模型的进展}

\textbf{MagicLens}是Google DeepMind提出的一种自监督图像检索模型，支持开放式的文本指令\cite{zhang2024magiclens}。它关注于这样的问题：给定一张查询图片和一段文本指令来描述想找的关系，检索出符合该关系的目标图片。为获取训练数据，研究者利用互联网中同一网页自然共现的图像对，假设这些图像对存在某种隐含关系。MagicLens在约3670万此类三元组上训练，从而学习在视觉特征之外捕捉指令语义所描述的深层关系。实验表明，MagicLens在包含多样语义关系的图像检索任务上取得了优异效果，而模型规模仅为SOTA方法的1/50，参数效率极高。

\textbf{FG-CLIP：}Fine-Grained CLIP是清华与奇虎360等提出的细粒度视觉-文本对齐模型\cite{xie2025fgclip}。它在OpenAI CLIP的基础上作了三点关键改进。首先，针对CLIP文本编码器只能处理最多77个词的问题，FG-CLIP扩展了位置编码长度至248，实现了长文本支持。其次，FG-CLIP引入大规模区域级别标注数据集FineHARD，通过在第二阶段训练中加入区域-文本对齐的对比学习，模型学会了将图像的局部区域与描述该区域的细粒度文本相匹配。第三，FG-CLIP设计了难负样本挖掘机制，生成语义上接近但不匹配的图文作为训练的细粒度负样本。实验结果显示，FG-CLIP在细粒度理解、开放词汇目标检测、图文检索等多项任务上全面超越原版CLIP。

\textbf{Long-CLIP：}Long-CLIP聚焦于提升CLIP对长文本的处理能力\cite{zhang2024longclip}。OpenAI原版CLIP由于位置编码限制，只能输入不超过77字的文本。Long-CLIP通过增加位置编码长度至248并利用含长描述的图文对数据进行再训练，使模型能够充分消化长文本信息。据报道，Long-CLIP将CLIP在长描述图文检索任务的R@5指标提升了约20\%，在传统短文本检索上也有约6\%的提升。

\textbf{EI-CLIP：}EI-CLIP提出了实体感知的干预式对比学习方法，用于电商跨模态检索任务\cite{ma2022eiclip}。该模型从语言学角度定义了实体感知检索任务，通过引入因果模型来增强跨模态检索的细粒度理解能力。

\subsection*{3.3 基于扩散模型的多模态数据生成与训练方法}

扩散模型（Diffusion Model）近年成为生成式多模态AI领域的焦点技术。其核心思想是通过逐步添加噪声然后反向去噪来从随机噪声生成逼真的数据，在图像生成等任务上表现卓越。

\textbf{文本-图像扩散模型：}代表性工作如OpenAI的GLIDE首次展示了利用扩散模型进行文本引导的图像生成与编辑的可能\cite{nichol2022glide}。Stable Diffusion利用预训练的VQGAN先将图像映射到低维语义潜空间，再在该潜空间执行扩散去噪，大幅降低了计算成本\cite{rombach2022stable}。自2022年开源发布以来，Stable Diffusion被广泛应用于文本到图像的生成。Google的PaLI模型则展示了大规模多模态预训练在视觉-语言任务上的强大能力\cite{chen2022pali}。

\textbf{条件扩散与控制生成：}ControlNet等架构进一步拓展了条件种类，允许使用边缘检测图、分割掩码、深度图甚至用户草图等作为条件，指导Stable Diffusion按照指定结构生成图像，实现高度可控的图像重构/生成\cite{zhang2023controlnet}。

\textbf{多模态伪样本生成与训练：}扩散模型的强大生成能力还被用来合成训练数据，以弥补真实多模态数据的不足。StableRep研究证明，使用Stable Diffusion合成的图像来训练模型，其效果可以媲美甚至超越使用真实图像训练\cite{tian2023stablerep}。这表明，在数据匮乏或标注昂贵的场景下，利用扩散模型大规模合成带标注的多模态数据是可行且有效的。

综上，多模态检索增强生成作为应对LLM幻觉和知识时效性的一项重要进展，正受到国内外广泛关注。在这一框架下，检索端有MagicLens、FG-CLIP、Long-CLIP等模型不断改进跨模态检索的精准度和丰富度；生成端有扩散模型提供强大的多模态生成与数据合成功能，为系统注入新鲜"素材"\cite{mei2025multimodalrag,abootorabi2025ask}。

\section*{四、本课题的研究内容}

本课题围绕多模态检索增强生成（mRAG）系统中的检索器性能这一关键瓶颈问题展开，其核心研究内容是：在多模态问答做题的场景下，系统性评估并提升"基于图像证据的检索器"相对于文本检索器与随机检索的有效性与可靠性。

\subsection*{4.1 基于 MRAG-Bench 的多模态检索评测任务建模}

MRAG-Bench数据集包含16,130张图像与1,353个人工标注的多选问答样本，覆盖9种典型多模态推理与知识依赖场景\cite{hu2024mragbench}。其核心特点在于：（1）每一道问题都经过人工设计，明确区分"是否需要视觉证据"；（2）多数问题在仅使用文本知识时存在歧义或错误风险；（3）图像作为外部知识证据对回答正确性具有决定性作用。因此，MRAG-Bench非常适合作为mRAG中"检索器是否真的找对证据"的评测基准，而不仅仅是测试最终模型能否"蒙对答案"。

在本课题中，MRAG-Bench被重构为如下评测流程：首先，检索器从图像池中检索Top-K候选图像；其次，将问题文本与检索到的图像拼接输入视觉-语言模型（VLM）；最后，由VLM输出多选题答案。其中，检索阶段是唯一变量，其余生成模型与提示词保持一致，从而将性能差异尽可能归因到"检索器质量"。

\subsection*{4.2 MagicLens 作为多模态检索器的引入与定位}

MagicLens是一种支持开放文本指令的图像检索模型\cite{zhang2024magiclens}，其核心能力并非单纯的视觉相似度匹配，而是学习"给定查询图像+文本指令时，检索与该语义关系最相关的目标图像"。这一能力与mRAG场景高度契合，原因在于：（1）mRAG的检索目标不是"最像的图"，而是"最有用的证据"；（2）多数MRAG-Bench问题隐含着关系性指令（如因果、属性、对比、背景）；（3）MagicLens的训练方式天然适配"关系驱动的证据检索"。

在本课题中，MagicLens被用于：（1）将文本和相关图像转化为检索指令；（2）从MRAG-Bench的图像池中检索Top-K候选图像；（3）作为mRAG系统中"图像检索模块"的核心实现。

\subsection*{4.3 检索器对比实验设计}

为验证MagicLens在mRAG中的实际价值，本课题设计了四类对照检索策略：

\textbf{（1）GT Oracle检索（上界）}：使用MRAG-Bench提供的人工标注正确图像作为"理论最优检索结果"，用于衡量下游模型在证据完美时的性能上限。

\textbf{（2）Random检索（下界）}：从图像库中随机采样Top-K图像作为"无语义信息检索"的下界基线，用于验证：视觉证据本身是否重要。

\textbf{（3）Text-only RAG（传统基线）}：不检索图像，仅检索或输入文本信息，模拟传统RAG/纯文本LLM的能力，用于对比：多模态检索是否优于文本增强\cite{lewis2020rag,gao2023rag}。

\textbf{（4）MagicLens多模态检索（研究对象）}：基于文本指令/图像+指令进行图像检索，检索结果作为视觉证据输入VLM。与上述三种方法在同一任务、同一生成模型下对比\cite{zhang2024magiclens}。

评价指标包括：（1）最终多选题准确率（Accuracy）；（2）不同场景（9类）下的性能差异；（3）相对于Text-only的增益幅度；（4）接近GT Oracle的程度（证据利用率）。

通过该设计，可以回答本课题的核心问题：在多模态问答中，一个"更懂指令的检索器"是否真的能显著提升mRAG的有效性？

\section*{五、本课题的实行方案、进度及预期效果}

\subsection*{5.1 实行方案}

实施方案总体框架如下：首先构建统一的评测管线，固定生成模型与提示词模板；其次，分别接入四种检索策略（GT Oracle、Random、Text-only、MagicLens）；最后，在MRAG-Bench全量数据上运行并统计各项指标。

该方案的特点是：（1）检索器为唯一变量，避免系统耦合；（2）评测结果可复现、可扩展到其他检索模型；（3）适合作为后续mRAG/检索器训练研究的基础设施。

\subsection*{5.2 研究进度安排}

\textbf{阶段一：基线复现与环境搭建（已完成/近期完成）}
\begin{itemize}
\item 搭建MRAG-Bench评测代码与数据管线
\item 跑通Text-only与Random检索基线
\item 验证不同VLM在无检索/随机检索下的稳定性
\end{itemize}

\textbf{阶段二：MagicLens集成与检索实验（核心阶段）}
\begin{itemize}
\item 接入MagicLens图像检索模型\cite{zhang2024magiclens}
\item 构建文本指令$\rightarrow$图像检索流程
\item 在MRAG-Bench全量问题上进行Top-K检索
\item 与GT/Random/Text-only进行系统对比
\end{itemize}

\textbf{阶段三：结果分析与消融实验}
\begin{itemize}
\item 按9种场景分析检索器收益分布
\item 分析MagicLens在哪些场景显著优于文本
\item 识别失败模式（检索对但生成错/检索错但生成对）
\end{itemize}

\subsection*{5.3 预期效果与成果形式}

\textbf{预期研究结论：}
\begin{enumerate}
\item 证明多模态检索器显著优于纯文本RAG
\item 量化MagicLens在mRAG中的实际收益
\item 揭示"哪些类型的问题真正依赖视觉证据"
\end{enumerate}

\textbf{预期成果：}
\begin{enumerate}
\item 一套可复现的mRAG检索评测流程
\item MRAG-Bench上系统性的检索器对比结果\cite{hu2024mragbench}
\item 为后续检索器训练、扩散模型驱动检索增强、具身/BCI/类脑mRAG提供坚实的实验基础\cite{xie2024embodiedrag,caria2025bci,genesis2025}
\end{enumerate}

\bibliographystyle{unsrt}
\bibliography{ref_ordered}

\end{document}
